%! Characteristics of Exponential Functions — Unit 6, Lesson 6.0 — Lesson Plan
\documentclass[10pt]{article}
\usepackage{algebra2-article}
\usepackage{algebra2-boxes}
\usepackage{graphicx}   % -article does NOT load graphicx; needed for thumbnails/screenshots
\graphicspath{{images/}}
\newcommand{\TallMath}[1]{$\displaystyle #1\rule[-1.4em]{0pt}{3.2em}$}
\newcommand{\CourseName}{Algebra 2}
\newcommand{\MeetingLength}{55 minutes}
\newcommand{\UnitNumberName}{Unit 6: Exponential Functions \quad}
\newcommand{\LessonNumberName}{Lesson 6.0: Characteristics of Exponential Functions}

\begin{document}
\begin{center}
    {\Large\bfseries \CourseName \\
    {\small\bfseries \UnitNumberName \LessonNumberName}}
\end{center}

\begin{tcolorbox}[colback=forestbg, colframe=forest, boxrule=0.9pt, arc=2mm,
    left=2mm, right=2mm, top=1.5mm, bottom=1.5mm]
    \textbf{Primary Objective:} Students can recognize an exponential function $f(x)=ab^{x}$ from its
    equation, its graph, or a \emph{table} --- via the \emph{constant ratio} that distinguishes it
    from a linear table's constant difference and a quadratic table's constant second difference ---
    and can read every characteristic off the two parents $y=2^{x}$ and $y=\left(\frac12\right)^{x}$,
    taught side by side: \emph{domain} (all reals, handed back after Unit~5 took half of it away),
    \emph{range} $(0,\infty)$, the \emph{$y$-intercept} that is always $a$, the \emph{absence} of any
    $x$-intercept, \emph{increasing/decreasing} and \emph{positive/negative} intervals, the total
    absence of \emph{extrema} and turning points, the \emph{horizontal asymptote} $y=0$, and
    \emph{end behavior} with two different \emph{kinds} of end --- one arm running to $\infty$ and
    the other flattening onto the asymptote. Students also distinguish \emph{growth} from
    \emph{decay} by comparing $b$ to $1$, and explain the lesson's central idea: an asymptote is a
    boundary the range \emph{excludes}, not a value the function \emph{reaches}, which is why an
    exponential graph has a visible floor and yet no absolute minimum.
    \\[2pt]
    \textbf{Standards (2023 VA SOL):} \textbf{A2.F.2} --- investigate and analyze characteristics of
    exponential functions graphically and algebraically. This Lesson~0 \emph{introduces} the single
    row the \S3 spine marks new for Unit~6 --- \textbf{growth vs.\ decay and the constant
    multiplicative rate} --- and \emph{revisits} the horizontal asymptote (\textbf{●} in Unit~4) in
    its new role as a \textbf{range boundary}. It covers \textbf{A2.F.2a} (domain, range, zeros,
    intercepts), \textbf{A2.F.2c} (increasing/decreasing intervals), \textbf{A2.F.2d} (absolute
    maxima and minima --- here, establishing that there are \emph{none}), \textbf{A2.F.2f} (evaluate
    $f(x)$ and interpret it in context --- the medication and depreciation models),
    \textbf{A2.F.2g} (end behavior), and \textbf{A2.F.2h}, which names exponential functions
    explicitly for determining horizontal asymptotes. \textbf{A2.F.2b} (compare and contrast the
    characteristics of function families) drives both the growth-vs.-decay table and the table
    fingerprint test. \textbf{A2.F.1a} (distinguish between the graphs of the parent functions) is
    the lesson's backbone, and \textbf{A2.F.1e} (compare and contrast graphs, tables, and equations)
    is assessed in the Activity Tier~A table. \emph{Note:} Unit~6 has no \texttt{A2.EO} or
    \texttt{A2.EI} standard --- the whole unit rests on \texttt{A2.F} and \texttt{A2.ST.2}.
\end{tcolorbox}

\renewcommand{\arraystretch}{1.4}
\begin{skillbox}[Priority Ideas \& Skills]{goldbox}
\begin{tabularx}{\linewidth}{@{}X|X@{}}
\textbf{Priority Ideas \& Skills} & \textbf{Key Understandings (the ``why'')} \\[2pt]
\hline
\vspace{-6pt}
\begin{itemize}[leftmargin=1.1em, nosep, after=\vspace{2pt}]
  \item Recognize $f(x)=ab^{x}$ by the variable's position --- in the \textbf{exponent}, not the base --- and identify $a$ and $b$.
  \item Use the \textbf{constant ratio} to name the family from a table, against linear (constant difference) and quadratic (constant second difference).
  \item Classify \textbf{growth} ($b>1$) vs.\ \textbf{decay} ($0<b<1$) by comparing the base to $1$.
  \item Read \textbf{domain}, \textbf{range}, \textbf{intercepts}, \textbf{increasing/decreasing} and \textbf{positive/negative} intervals off both exponential parents.
  \item State the \textbf{horizontal asymptote} and say which \emph{side} the graph approaches it on.
  \item Give \textbf{end behavior} in both directions --- including the arm that flattens rather than running to $\pm\infty$.
  \item Explain why the family has \textbf{no zeros} and \textbf{no extrema}, ever.
  \item \textbf{Compare and contrast} growth with decay, and exponential with the families of Units 1--5.
\end{itemize}
&
\vspace{-6pt}
\begin{itemize}[leftmargin=1.1em, nosep, after=\vspace{2pt}]
  \item \emph{The base runs the show.} Whether the graph rises or falls, and which side hugs the asymptote, both trace back to $b$ versus $1$. $a$ only sets the starting height.
  \item Every family so far grew by \textbf{adding}. This one grows by \textbf{multiplying} --- that is the difference, not ``fast'' versus ``slow.''
  \item \textbf{An asymptote is a boundary, not a value.} That single sentence answers ``why $(0,\infty)$ and not $[0,\infty)$?'' and ``why is there no minimum?'' at the same time.
  \item Unit 5's endpoint \emph{was} on the graph, so it was an absolute minimum. Unit 6's floor is not on the graph, so it is not.
  \item Unit 5 took away half the \textbf{domain}; Unit 6 hands the domain back and takes half the \textbf{range}. The escalation reverses.
  \item ``End behavior'' has had two answers all year. Today one arm answers $\pm\infty$ and the other answers \emph{it flattens} --- the mirror of 5.0's ``the graph stops.''
  \item A positive base raised to \emph{any} real power is positive. That is the whole reason there is never an $x$-intercept.
\end{itemize}
\end{tabularx}
\end{skillbox}

\begin{skillbox}[{Vocabulary, Concepts, \& Theorems}]{sky}
\renewcommand{\arraystretch}{1.3}
\begin{tabularx}{\linewidth}{@{}l X@{}}
\textbf{Exponential function} & A function with the variable in the \emph{exponent}: $f(x)=ab^{x}$, with $a\ne0$, $b>0$, $b\ne1$. \\
\textbf{Base $b$ (growth/decay factor)} & The number being raised to the power --- what you multiply the output by each time $x$ increases by $1$. \\
\textbf{Initial value $a$} & The output at $x=0$. Since $b^{0}=1$ for every legal base, $f(0)=a$, so $a$ is always the $y$-intercept. \\
\textbf{Exponential growth} & $b>1$: each step multiplies by more than $1$, so outputs increase without bound. \\
\textbf{Exponential decay} & $0<b<1$: each step multiplies by a proper fraction, so outputs shrink toward $0$. \emph{Compare $b$ to $1$, not to $0$.} \\
\textbf{Constant ratio} & Dividing consecutive outputs (evenly spaced inputs) always gives the same number, $b$ --- the table fingerprint of the family. \\
\textbf{Horizontal asymptote} & A horizontal line the graph approaches but never reaches. For $y=ab^{x}$ it is always $y=0$, and it is exactly the edge the range excludes. \\
\textbf{Growth parent} & $y=2^{x}$: domain all reals, range $(0,\infty)$, $y$-int $(0,1)$, no zeros, increasing, no extrema, HA $y=0$ approached on the left. \\
\textbf{Decay parent} & $y=\left(\frac12\right)^{x}$: everything identical except decreasing, with the HA approached on the right. Also equals $2^{-x}$. \\
\textbf{Why no extrema} & A minimum or maximum must be a value the function \emph{attains}. The graph attains every value in $(0,\infty)$ and never attains $0$, so neither exists. \\
\end{tabularx}
\end{skillbox}

\begin{fixedskillbox}[Activate Prior Knowledge \& Spiral Review]{forestbg}
    The Warm-Up rehearses the three skills this lesson stands on, one per new idea: \textbf{(1)}
    evaluating $2^{3}$, $2^{0}$, $2^{-1}$, $2^{-2}$, $2^{-4}$, $2^{-8}$, and $5^{0}$ --- so students
    discover \emph{numerically} that a negative exponent produces a shrinking \emph{fraction} rather
    than a negative number, which is the entire asymptote story in miniature, while $2^{0}$ and
    $5^{0}$ plant the $y$-intercept; \textbf{(2)} comparing an ``add $3$'' table with a ``times $3$''
    table and naming the operation --- the fingerprint test, with the linear case reactivated from
    Lesson~1.0; and \textbf{(3)} a doubling-rumor table starting at $3$ people, which plants both
    $a$ (the value at hour $0$) and $b$ (the constant multiplier) without the notation. Debrief all
    three aloud, but do \emph{not} supply the words ``exponential,'' ``growth factor,'' or
    ``asymptote'' --- let the Hook graphs earn them.
\end{fixedskillbox}

\begin{skillbox}[Hook]{forestbg}
    Put the two parent graphs on the board side by side: $y=2^{x}$ and $y=\left(\frac12\right)^{x}$.
    Open by naming what has changed since yesterday's unit: \textbf{``Unit 5's square root graph gave
    up half the inputs --- it just stopped. Look at these two. They run all the way across, so the
    domain is back. What did they give up instead?''} Let students hunt. Push until someone says the
    graphs never go below the $x$-axis, then press: \textbf{``So how much of the $y$-axis are we
    losing?''} (All of it below zero --- half the outputs.) Now pull the reason out of the Warm-Up:
    \textbf{``Item 1 already told you. What was $2^{-8}$?''} ($\frac{1}{256}$ --- tiny, but positive.)
    \textbf{``Can you make it zero?''} (No --- halving a positive number forever never reaches zero.)
    Land the trade explicitly: \emph{Unit 5 took away half the domain; Unit 6 hands the domain back
    and takes half the range instead.} Then set up the day's real question by pointing at the flat
    left arm of $y=2^{x}$: \emph{that graph clearly has a floor --- so what is its minimum?} Take the
    wrong answer (``zero'') without correcting it yet. Section 4 of the notes is where it gets
    settled, and the lesson is more memorable if the class has already committed to an answer.
\end{skillbox}

\begin{skillbox}[Lesson]{forestbg}
\begin{multicols}{2}
\textbf{1. What makes a function exponential --- and the table fingerprint.} Contrast $y=x^{2}$ with
$y=2^{x}$ on the board: the variable has moved from the \textbf{base} to the \textbf{exponent}.
Introduce $f(x)=ab^{x}$ with $a\ne0$, $b>0$, $b\ne1$, and get $a$'s meaning immediately from
Warm-Up item 1: $b^{0}=1$ for \emph{every} legal base, so $f(0)=a$ and $a$ is always the
$y$-intercept. Keep the restrictions brief --- they are guardrails, not content ($1^{x}=1$ is a
horizontal line; $(-4)^{1/2}$ is not real). Then spend the time on the \textbf{fingerprint table}:
linear keeps a constant \emph{difference} (Unit 1), quadratic a constant \emph{second} difference
(Unit 2), exponential a constant \textbf{ratio}. Insist on the word \emph{multiply}. Close with the
growth/decay rule --- \textbf{compare $b$ to $1$, never to $0$}.

\textbf{2. The growth parent.} Build $y=2^{x}$ from the table $(-3,\frac18)$, $(-1,\frac12)$,
$(0,1)$, $(2,4)$, $(3,8)$ and read every feature: domain all reals, range $(0,\infty)$ with the
round bracket flagged out loud, $y$-int $(0,1)$, \textbf{no} $x$-intercept, increasing everywhere,
\textbf{no extrema at all}, HA $y=0$. Spend real time on end behavior: as $x\to+\infty$ it climbs
forever and \emph{fast}; as $x\to-\infty$ it does \emph{not} run down --- it runs \emph{flat}. Say
plainly that ``$-\infty$'' is the reflex answer from Unit 3 and it is wrong here.

\columnbreak

\textbf{3. The decay parent.} Same treatment for $y=\left(\frac12\right)^{x}$, and let students
notice how few answers change: only the direction and which side hugs the asymptote. Close the
section by rewriting $\left(\frac12\right)^{x}=\frac{1}{2^{x}}=2^{-x}$ --- a Unit 5 exponent
identity, not a new rule --- and note that the two parents are one graph reflected. Flag Lesson 6.2
and move on; the reflection \emph{rule} is not today's business.

\textbf{4. The asymptote --- the heart of the lesson.} Put $2^{-1}, 2^{-4}, 2^{-10}, 2^{-20}$ on the
board and let the class watch the outputs shrink without ever turning zero or negative. Then settle
the Hook's open question: the range is $(0,\infty)$, \emph{not} $[0,\infty)$, and there is
\textbf{no absolute minimum}, because a minimum must be a value the function \emph{reaches}. Draw
Unit 5's $\sqrt{x}$ beside $2^{x}$: solid endpoint versus dashed asymptote. Add the Unit 4 contrast
--- a rational graph may cross its HA in the middle; $y=ab^{x}$ never crosses its own.

\textbf{5. The full feature list and the contrast.} Read every characteristic off the anchor
$f(x)=4\left(\frac12\right)^{x}$ (decay, $y$-int $(0,4)$, range $(0,\infty)$, decreasing, HA $y=0$
approached on the right), then pull $b$ off the \emph{graph} by dividing consecutive outputs ---
the skill 6.1 and 6.5 both need. Finish with the \textbf{A2.F.2b} growth-vs-decay table, making the
point that most rows come out identical, and one sentence against the earlier units: a polynomial
turned around, a rational function broke apart, a radical function began and stopped --- an
exponential does none of it. Preview Lesson 6.1, \emph{exponential growth \& decay}.
\end{multicols}
\end{skillbox}

\begin{skillbox}[Active Monitoring]{redbox}
Circulate during the notes practice and Tier work. Watch for and correct:
\begin{itemize}[leftmargin=1.2em, nosep]
  \item writing the range as $[0,\infty)$ --- treating the asymptote as an attained value; this is \emph{the} signature error of the lesson, and it is the Unit 5 endpoint habit carried over unexamined;
  \item claiming an absolute minimum of $0$ (same error, different costume) --- ask ``what input gives you exactly zero?'';
  \item answering ``$y\to-\infty$'' for the flattening arm --- the reflex from odd-degree polynomials; the left arm of a growth graph runs \emph{flat}, not \emph{down};
  \item deciding growth vs.\ decay from $a$ instead of $b$ (e.g.\ calling $y=0.4(3)^{x}$ decay because $0.4$ is small), or comparing $b$ to $0$ instead of to $1$;
  \item hunting for an $x$-intercept and being unwilling to write ``none'';
  \item reading a negative exponent as a negative output --- send them back to Warm-Up item 1;
  \item calling a fast-growing table exponential without checking the ratios (Homework item 4, Table C is quadratic);
  \item naming the base of $y=-2\cdot3^{x}$ as $-2$ --- the negative belongs to $a$, and it decides which \emph{side} of the asymptote the graph lives on, not the base.
\end{itemize}
Cold-call prompts: ``Is $b$ bigger or smaller than $1$ --- and so which is it?'' ``Which side does
this graph approach its asymptote on, and how do you know from the equation?'' ``What input makes
this function output zero?'' ``Say what happens on the left, exactly.'' ``Difference or ratio ---
which one is constant here?''
\end{skillbox}

\begin{skillbox}[Group Work \& Differentiation]{redbox}
\textbf{Reading Exponential Graphs} activity (tiered; mirrors the notes).
\begin{multicols}{3}
\textbf{Tier R --- Remediate.} From one pre-drawn graph of $g(x)=3^{x}$: identify $a$ and $b$ and
classify growth vs.\ decay, give domain and range, check the $y$-intercept against $a$, then
\emph{attempt} an $x$-intercept by evaluating $3^{-1}$ and $3^{-4}$ and conclude there is
\textbf{none}. Closes on the asymptote and on end behavior in both directions --- with the left one
required to be answered ``it flattens onto $y=0$,'' not ``$-\infty$.''

\columnbreak
\textbf{Tier A --- Approaching.} An eleven-row feature table for $p(x)=4^{x}$ beside
$q(x)=8\left(\frac14\right)^{x}$, chosen so that \emph{five rows come out identical}. Then two
\emph{justifications}: point to the feature of each equation that decides which one rises (the base
--- and explicitly \emph{not} $a$, which is deliberately larger for the falling one), and list the
five shared rows, explaining what they establish about every $y=ab^{x}$ with $a>0$. Justify 2 is the
real \textbf{A2.F.2b} assessment of the day.

\columnbreak
\textbf{Tier E --- Extension.} Equation-only work: a four-row table (base / $b$ vs.\ $1$ /
$y$-intercept / range) over $7(1.5)^{x}$, $7(0.5)^{x}$, $-2\cdot3^{x}$, and $\left(\frac54\right)^{x}$,
plus ``which one's range is not $(0,\infty)$, and how do you know before graphing?'' (the sign of
$a$). Then a medication model $A(t)=200(0.7)^{t}$: identify and interpret $a$ and $b$, confirm the
constant ratio, evaluate and interpret $A(3)$, contrast the algebraic domain with the contextual
one, say what the asymptote claims about the drug and whether it is believable, and refute a
classmate's ``it loses 60 mg an hour'' with the shrinking differences.
\end{multicols}
\end{skillbox}

\begin{skillbox}[Individual Work \& Assessment]{redbox}
\textbf{Exit Ticket} (independent, no notes): from one graphed exponential
$f(x)=8\left(\frac12\right)^{x}$, give $a$, $b$, and growth vs.\ decay; domain, range, and the
horizontal asymptote; both intercepts (one of which is \emph{none}); increasing or decreasing; and
end behavior in both directions. Then \textbf{explain} why $f$ has no $x$-intercept \emph{and} why
there is no absolute minimum despite the visible floor. Then an \textbf{SOL-style multiple-choice}
item on the range of $g(x)=6\left(\frac13\right)^{x}$ (answer $(0,\infty)$) whose distractors each
break exactly one idea: $[0,\infty)$ treats the asymptote as attained, $(-\infty,\infty)$ forgets
the floor and answers with the domain, and $(-\infty,0)$ confuses ``decay'' with ``negative.''
Collect and sort into ``got it / closed the bracket / decay-means-negative / hunted for an
$x$-intercept'' to decide whether Lesson~6.1 opens with a re-teach. Also check that students pulled
$a=8$ off the $y$-intercept --- 6.1 asks them to \emph{build} $y=ab^{x}$ from scratch, and anyone who
cannot find $a$ on a graph will stall immediately.
\end{skillbox}

\begin{skillbox}[Reinforcement \& Extension]{goldbox}
\textbf{Homework:} a four-row equation-only table (base / growth or decay / $y$-intercept / which
side approaches the asymptote) over $6\cdot2^{x}$, $6\left(\frac12\right)^{x}$,
$\left(\frac45\right)^{x}$, and $0.4(3)^{x}$ --- the last one deliberately pairing a small $a$ with a
large $b$ --- followed by the observation that all four share a domain, a range, and an asymptote; a
full feature read off the growth graph $m(x)=4\cdot2^{x}$; a matched read off the decay graph
$n(x)=9\left(\frac13\right)^{x}$, closing with which side each of the two approaches $y=0$ on and
what in the equation decides it; a three-table \textbf{fingerprint} item (linear / exponential /
quadratic, with the quadratic disguised as fast growth); and one justification item on why Unit 5's
endpoint produced an absolute minimum while Unit 6's asymptote does not --- the unit hinge, and the
item to grade closely. \textbf{Extension:} interpret a depreciation model $V(t)=1200(0.75)^{t}$ ---
$a$ and $b$ and why $b$ is $0.75$ rather than $0.25$, the constant ratio, evaluating and
interpreting $V(2)$, refuting ``it loses \$300 a year'' with the shrinking differences, what the
asymptote claims about a laptop that is never worth nothing, and the algebraic domain versus the
contextual one. \textbf{Preview:} Lesson~6.1, \emph{Exponential Growth \& Decay} --- building
$y=ab^{x}$ from a table, from two points, and from a story, where the whole lesson turns on one
translation: ``up $8\%$'' means $b=1.08$, not $b=0.08$, and ``down $15\%$'' means $b=0.85$, not
$b=1.15$. Today's medication model already did it once at $70\%$ remaining.
\end{skillbox}
\end{document}
